\chapter*{Abstract}

Sign language recognition is an important accessibility challenge, as it can help reduce communication barriers between deaf or hard-of-hearing signers and non-signing people. This project develops and evaluates a research prototype for word-level American Sign Language recognition from short video clips. The
final approach extracts 21 three-dimensional hand landmarks per frame with MediaPipe Hands, normalises the resulting sequences, and classifies ASL glosses using a bidirectional GRU with a Temporal Attention mechanism. The work also documents the methodological evolution from YOLO-based hand detection and static
letter recognition toward landmark-based temporal modelling, where the main challenge lies in representing sign motion rather than localising hands. The system was validated through experiments on WLASL and MS-ASL, including architecture comparison, sequence-length analysis, cross-dataset evaluation, and
vocabulary scaling. Results indicate promising performance for reduced vocabularies and reproducible behaviour across datasets, but also clear limitations as the number of classes increases. The landmark representation provides a privacy-friendly abstraction, although the prototype remains a research system
rather than a production-ready translator.

\begin{keywords}{Keywords}
  \mykeywords{}
\end{keywords}
